% Part: lambda-calculus
% Chapter: introduction
% Section: lambda-definability

\documentclass[../../../include/open-logic-section]{subfiles}

\begin{document}

\olfileid[hi]{lam}{int}{rep}
\olsection{\usetoken{S}{lambda definable} अंकगणितीय फलन}

लैम्ब्डा कलन संगणन के मॉडल का काम कैसे कर सकता है? आरंभ में तो यह भी स्पष्ट
नहीं कि इस कथन का अर्थ कैसे समझें। प्राकृतिक संख्याओं पर संगणनीयता की बात
करने के लिए हमें उन संख्याओं का उपयुक्त निरूपण खोजना होगा। यह एक ऐसा
निरूपण है जो आश्चर्यजनक रूप से अच्छा काम करता है।

\begin{defn}
प्रत्येक प्राकृतिक संख्या~$n$ के लिए \emph{चर्च अंक}
$\num{n}$ को लैम्ब्डा पद $\lambd[x][\lambd[y][(x(x(x(\dots x(y)))))]]$
परिभाषित करें, जहाँ कुल $n$ बार $x$ आता है।
\end{defn}

पद $\num{n}$ ``पुनरावर्तक'' हैं: निवेश $f$ पर $\num{n}$ वह फलन लौटाता है
जो $y$ को $f^n(y)$ पर प्रतिचित्रित करता है। ध्यान दें कि प्रत्येक अंक
प्रसामान्य है। अब हम बता सकते हैं कि किसी लैम्ब्डा पद द्वारा प्राकृतिक
संख्याओं पर कोई फलन ``संगणित करने'' का क्या अर्थ है।

\begin{defn}
$f(x_0, \dots, x_{k-1})$, ~$\Nat$ से $\Nat$ में कोई $n$-स्थानी आंशिक फलन
हो। हम कहते हैं कि $\lambd$-पद~$F$, ~$f$ को \emph{!!{lambda define}}
करता है, तभी और केवल तभी जब प्राकृतिक संख्याओं के प्रत्येक अनुक्रम $n_0$,
\dots,~$n_{k-1}$ के लिए,
\[
F\, \num{n_0}\, \num{n_1} \dots \num{n_{k-1}} \red \num{f(n_0, n_1, \dots,
  n_{k-1})}
\]
यदि $f(n_0, \dots, n_{k-1})$ परिभाषित है; और अन्यथा $F, \num{n_0}\,
\num{n_1} \dots \num{n_{k-1}}$ का कोई प्रसामान्य रूप न हो।
\end{defn}

\begin{thm}
\ollabel{thm:lambda-def}
कोई फलन $f$ आंशिक संगणनीय फलन तभी और केवल तभी है जब वह किसी लैम्ब्डा पद
द्वारा !!{lambda defined} हो।
\end{thm}

\begin{explain}
यह प्रमेय कुछ चकित करने वाला है। संगणन के मॉडल के रूप में लैम्ब्डा कलन
अत्यंत सरल कलन है; इसमें केवल लैम्ब्डा अमूर्तन और अनुप्रयोग संक्रियाएँ हैं!{}
फिर भी इन सीमित साधनों से किसी भी संगणनात्मक प्रक्रिया को कार्यान्वित करना
संभव है।
\end{explain}

\end{document}
